\section{Resultados}
En esta secci\'on se presentan los resultados obtenidos a partir de seis escenarios experimentales ejecutados sobre la plataforma ThingsBoard conectada a la red IEEE~802.11ah (HaLow). Cada prueba explora un par\'ametro distinto de carga o configuraci\'on ---n\'umero de dispositivos, intervalo de env\'io, tama\~no del \emph{payload}, l\'imites de tasa y concurrencia as\'incrona--- con el objetivo de delimitar la regi\'on de operaci\'on segura del sistema sin degradar la calidad de servicio percibida por los dispositivos simulados.

En todos los casos se graficaron las m\'etricas principales (latencia promedio, percentil 95, \emph{throughput} y estabilidad de la red HaLow) y se consolidaron res\'umenes num\'ericos a partir de los archivos CSV producidos por los scripts en \texttt{data/metrics/}.

\subsection{Prueba 1: l\'inea base con 20 dispositivos}
La primera prueba corresponde a una l\'inea base con \num{20} dispositivos simulados, cada uno publicando telemetr\'ia peri\'odica hacia ThingsBoard a trav\'es del broker MQTT expuesto sobre la red HaLow. Esta corrida permite verificar que la instrumentaci\'on funciona correctamente y que, bajo baja concurrencia, el sistema opera sin p\'erdida de mensajes ni latencias excesivas.

La \Cref{tab:prueba1} resume los resultados agregados. Se observa que los \num{20} clientes se mantienen conectados durante aproximadamente \num{162}~\si{\second}, con \num{300} mensajes exitosos y sin publicaciones fallidas. La latencia promedio se sit\'ua en torno a \SI{6.45}{\milli\second}, con un percentil 95 por debajo de \SI{10}{\milli\second}, mientras que el \emph{throughput} agregado es cercano a \SI{1.7}{\text{msg/s}}. La capa f\'isica HaLow presenta un RSSI promedio de aproximadamente \SI{-24}{\decibelm} y niveles de ruido alrededor de \SI{-81}{\decibelm}, lo que indica un enlace c\'omodo lejos del umbral de sensibilidad.

\begin{table}[!t]
\centering
\caption{Resumen de la prueba 1 (l\'inea base con 20 dispositivos).}
\label{tab:prueba1}
\begin{tabular}{@{}lrrrrrr@{}}
\toprule
Configuraci\'on & Dispositivos & Duraci\'on (s) & Mensajes exitosos & Latencia prom. (ms) & Throughput (msg/s) \\
\midrule
n20 & 20 & 162.38 & 300 & 6.45 & 1.70 \\
\bottomrule
\end{tabular}
\end{table}

\subsection{Prueba 2: escalamiento del n\'umero de dispositivos}
La segunda prueba analiza el impacto de incrementar el n\'umero de dispositivos concurrentes manteniendo constantes el intervalo de env\'io y el tama\~no del \emph{payload}. Se consideran cuatro configuraciones: \num{50}, \num{100}, \num{200} y \num{400} dispositivos activos en paralelo.

En todos los casos la duraci\'on de la simulaci\'on se mantiene cercana a \num{192}~\si{\second}. La \Cref{tab:prueba2} presenta los resultados agregados. Se aprecia un crecimiento casi lineal del \emph{throughput} ---de aproximadamente \SI{8.6}{\text{msg/s}} con \num{50} clientes hasta alrededor de \SI{69}{\text{msg/s}} con \num{400} dispositivos--- sin p\'erdida de mensajes ni desconexiones de clientes.

La latencia promedio se mantiene por debajo de \SI{5}{\milli\second} hasta \num{100} dispositivos, aumenta hasta unos \SI{11.4}{\milli\second} en la configuraci\'on de \num{200} clientes y desciende nuevamente a alrededor de \SI{6.1}{\milli\second} con \num{400} dispositivos. El percentil 95 de latencia alcanza un m\'aximo de aproximadamente \SI{37}{\milli\second} en la configuraci\'on intermedia de \num{200} nodos, pero sin superar en ning\'un caso valores que comprometan el uso de ThingsBoard para telemetr\'ia peri\'odica.

Desde la perspectiva de la red HaLow, el RSSI promedio oscila en torno a \SI{-32}{\decibelm} con niveles de ruido cercanos a \SI{-81}{\decibelm}, y las tasas de transmisi\'on se mantienen alrededor de \SI{14.5}{\mega\bit\per\second}, sin evidencias de saturaci\'on en la capa de enlace.

\begin{table}[!t]
\centering
\caption{Prueba 2: escalamiento del n\'umero de dispositivos conectados.}
\label{tab:prueba2}
\begin{tabular}{@{}lrrrrrr@{}}
\toprule
Configuraci\'on & Dispositivos & Duraci\'on (s) & Mensajes exitosos & Latencia prom. (ms) & Throughput (msg/s) \\
\midrule
n50  &  50 & 192.38 &  1800 &  4.91 &  8.64 \\
n100 & 100 & 192.36 &  3600 &  4.83 & 17.28 \\
n200 & 200 & 192.36 &  7200 & 11.38 & 34.55 \\
n400 & 400 & 192.36 & 14400 &  6.08 & 69.11 \\
\bottomrule
\end{tabular}
\end{table}

\subsection{Prueba 3: variaci\'on del intervalo de env\'io}
En la tercera prueba se fija el n\'umero de dispositivos en \num{200} y se var\'ia el intervalo de env\'io de telemetr\'ia. Se eval\'uan cuatro intervalos: \SI{10}{\second}, \SI{5}{\second}, \SI{2}{\second} y \SI{1}{\second} entre publicaciones.

La \Cref{tab:prueba3} resume los resultados. Al aumentar la frecuencia de env\'io (intervalos m\'as cortos), el \emph{throughput} agregado crece desde alrededor de \SI{17}{\text{msg/s}} con \SI{10}{\second} hasta aproximadamente \SI{172}{\text{msg/s}} con \SI{1}{\second}. Este incremento viene acompa\~nado de cambios en la latencia: con intervalos de \SI{10}{\second} la latencia promedio se sit\'ua cerca de \SI{5.6}{\milli\second}, mientras que para \SI{5}{\second} y \SI{2}{\second} aumenta hasta unos \SI{8.7}{\milli\second} y \SI{9.7}{\milli\second}, respectivamente. Para el intervalo de \SI{1}{\second} la latencia promedio vuelve a valores pr\'oximos a \SI{5.9}{\milli\second}, aunque el percentil 95 puede superar los \SI{14}{\milli\second} en los escenarios m\'as exigentes.

En todos los casos no se registran publicaciones fallidas ni dispositivos ca\'idos, lo que sugiere que, para \num{200} nodos, el l\'imite de carga est\'a m\'as condicionado por el \emph{back-end} de aplicaci\'on y el procesamiento interno de ThingsBoard que por la infraestructura de red HaLow, la cual mantiene tasas de transmisi\'on en torno a \SI{14}{\mega\bit\per\second} con niveles de RSSI estables.

\begin{table}[!t]
\centering
\caption{Prueba 3: variaci\'on del intervalo de env\'io con 200 dispositivos.}
\label{tab:prueba3}
\begin{tabular}{@{}lrrrrrr@{}}
\toprule
Intervalo & Dispositivos & Duraci\'on (s) & Mensajes exitosos & Latencia prom. (ms) & Throughput (msg/s) \\
\midrule
int10 & 200 & 192.35 &  3600 &  5.56 &  17.28 \\
int5  & 200 & 192.38 &  7200 &  8.67 &  34.54 \\
int2  & 200 & 192.39 & 18000 &  9.67 &  86.35 \\
int1  & 200 & 192.38 & 35800 &  5.91 & 172.10 \\
\bottomrule
\end{tabular}
\end{table}

\subsection{Prueba 4: tama\~no del \emph{payload}}
La cuarta prueba estudia el efecto del tama\~no del mensaje de telemetr\'ia sobre el comportamiento del sistema. Se mantienen \num{200} dispositivos y un intervalo de env\'io fijo, pero se var\'ia el \emph{payload} entre tres perfiles: \textit{small}, \textit{medium} y \textit{large}.

Los resultados de la \Cref{tab:prueba4} muestran que el n\'umero de mensajes exitosos y la duraci\'on de las corridas se mantienen pr\'acticamente constantes (alrededor de \num{7200} mensajes en unos \num{192}~\si{\second}), mientras que la latencia s\'i se ve afectada. En el perfil \textit{small} la latencia promedio se sit\'ua en torno a \SI{10.4}{\milli\second}, en \textit{medium} aumenta hasta aproximadamente \SI{14.1}{\milli\second} con un percentil 95 cercano a \SI{54}{\milli\second}, y en \textit{large} desciende nuevamente hasta unos \SI{8.5}{\milli\second} con p95 por debajo de \SI{21}{\milli\second}.

Este comportamiento sugiere que el coste computacional asociado al procesamiento de mensajes de tama\~no intermedio puede ser m\'as relevante que el mero volumen de bytes transmitidos, al menos en la configuraci\'on espec\'ifica de ThingsBoard y del simulador empleada.

\begin{table}[!t]
\centering
\caption{Prueba 4: variaci\'on del tama\~no de \emph{payload} con 200 dispositivos.}
\label{tab:prueba4}
\begin{tabular}{@{}lrrrrrr@{}}
\toprule
Tama\~no & Dispositivos & Duraci\'on (s) & Mensajes exitosos & Latencia prom. (ms) & Throughput (msg/s) \\
\midrule
small  & 200 & 192.39 & 7200 & 10.40 & 34.54 \\
medium & 200 & 192.39 & 7200 & 14.06 & 34.54 \\
large  & 200 & 192.39 & 7200 &  8.47 & 34.54 \\
\bottomrule
\end{tabular}
\end{table}

\subsection{Prueba 5: l\'imite de tasa y agregaci\'on}
En la quinta prueba se comparan dos estrategias de env\'io con \num{200} dispositivos: una configuraci\'on sin agregaci\'on ni l\'imite expl\'icito de tasa (\textit{sin\_aggregate}) y otra en la que se introducen restricciones de \emph{rate limit} a nivel de \emph{broker} o aplicaci\'on (\textit{con\_rate\_limit}).

Como se aprecia en la \Cref{tab:prueba5}, el escenario sin l\'imite expl\'icito alcanza cerca de \num{9600} mensajes exitosos en unos \num{252}~\si{\second}, con un \emph{throughput} promedio de alrededor de \SI{35}{\text{msg/s}} y una latencia media de \SI{8.2}{\milli\second}. Al activar el \emph{rate limit}, el n\'umero total de mensajes se reduce a aproximadamente \num{4800} y el \emph{throughput} promedio se sit\'ua en torno a \SI{17.7}{\text{msg/s}}, mientras que la latencia aumenta ligeramente hasta unos \SI{10}{\milli\second}.

Desde el punto de vista de calidad de servicio, ambos modos mantienen p\'erdida de mensajes nula; sin embargo, el l\'imite de tasa introduce un compromiso claro entre volumen de telemetr\'ia transmitido y latencia, lo que resulta relevante si se desea proteger al \emph{back-end} frente a picos de carga a costa de reducir la resoluci\'on temporal de las mediciones.

\begin{table}[!t]
\centering
\caption{Prueba 5: comparaci\'on entre env\'io directo y l\'imite de tasa (\textit{rate limit}).}
\label{tab:prueba5}
\begin{tabular}{@{}lrrrrrr@{}}
\toprule
Modo & Dispositivos & Duraci\'on (s) & Mensajes exitosos & Latencia prom. (ms) & Throughput (msg/s) \\
\midrule
sin\_aggregate  & 200 & 252.39 & 9600 &  8.17 & 35.38 \\
con\_rate\_limit & 200 & 252.40 & 4800 & 10.06 & 17.69 \\
\bottomrule
\end{tabular}
\end{table}

\subsection{Prueba 6: escenario as\'incrono con m\'ultiples lotes}
La \'ultima prueba eval\'ua un escenario m\'as cercano a una situaci\'on real, en el que grupos de dispositivos se activan en momentos distintos. Se consideran tres configuraciones: \texttt{async\_n200\_s800}, con \num{200} nodos y un desplazamiento temporal, y dos variantes con \num{400} dispositivos (\texttt{async\_n400\_s0} y \texttt{async\_n400\_s400}) que representan diferentes patrones de arranque.

Seg\'un la \Cref{tab:prueba6}, la configuraci\'on de \num{200} dispositivos alcanza del orden de \num{12000} mensajes exitosos en \num{180}~\si{\second}, con un \emph{throughput} medio de \SI{66.7}{\text{msg/s}} y latencia promedio cercana a \SI{6.1}{\milli\second}. Las configuraciones con \num{400} dispositivos duplican el \emph{throughput} hasta aproximadamente \SI{133}{\text{msg/s}}, con latencias promedio entre \SI{8.3}{\milli\second} y \SI{8.6}{\milli\second} y percentiles 95 en torno a \SI{25}{\milli\second}.

En todos los casos la p\'erdida de mensajes permanece en cero y no se observan dispositivos marcados como fallidos. La red HaLow opera con un RSSI promedio pr\'oximo a \SI{-21.5}{\decibelm} y tasas de transmisi\'on superiores a \SI{15}{\mega\bit\per\second}, por lo que el cuello de botella del sistema se asocia nuevamente al procesamiento de mensajes en el lado de ThingsBoard y del generador de carga, m\'as que a la capacidad del enlace inal\'ambrico.

\begin{table}[!t]
\centering
\caption{Prueba 6: escenario as\'incrono con m\'ultiples lotes de dispositivos.}
\label{tab:prueba6}
\begin{tabular}{@{}lrrrrrr@{}}
\toprule
Configuraci\'on & Dispositivos & Duraci\'on (s) & Mensajes exitosos & Latencia prom. (ms) & Throughput (msg/s) \\
\midrule
async\_n200\_s800  & 200 & 180.03 & 12000 &  6.09 & 66.67 \\
async\_n400\_s0    & 400 & 180.03 & 24000 &  8.36 & 133.26 \\
async\_n400\_s400  & 400 & 180.03 & 24000 &  8.55 & 133.32 \\
\bottomrule
\end{tabular}
\end{table}
